% TOP_PROBLEM: {"id": 3433, "title": "Jacob Palis Conjecture(Finitude of Attractors)(Dynamical Systems)", "category_id": 7, "category": {"id": 7, "name": "topology", "display_name": "Topology", "description": "Properties preserved under continuous deformations.", "slug": "topology", "order_index": 7, "created_at": "2026-07-31T15:26:25.671Z"}, "status": "open"}
% FIRST_POSTED: 2026-09-27
% Standalone research entry 205; batch 203--207.
% Original section material: Pasted text(3).txt, source edition 22 September 2026.
% Research additions: 27 September 2026.
% Compile twice with: lualatex 3433.tex
% !TeX program = lualatex
% Compile twice: lualatex open_problems_500.tex
% All references and prose are embedded; no BibTeX database or external inputs.
\documentclass[11pt,a4paper]{article}
\usepackage[margin=24mm,headheight=14pt]{geometry}
\usepackage{fontspec}
\setmainfont{Latin Modern Roman}
\setsansfont{Latin Modern Sans}
\setmonofont{Latin Modern Mono}
\usepackage{amsmath,amssymb,mathtools}
\usepackage{unicode-math}
\setmathfont{Latin Modern Math}
\usepackage{microtype}
\usepackage{enumitem,xurl,needspace}
\usepackage[dvipsnames]{xcolor}
\usepackage{hyperref}
\hypersetup{unicode=true,colorlinks=true,linkcolor=MidnightBlue,urlcolor=MidnightBlue,
 pdftitle={500 Mathematical Problem Statements},pdfauthor={Source-annotated compilation},bookmarksnumbered=true}
\usepackage{bookmark}
\usepackage{tocloft}
\setlength{\cftsecnumwidth}{3em}
\cftsetpnumwidth{2.5em}
\cftsetrmarg{4em}
\usepackage{titlesec}
\titleformat{\section}{\Large\bfseries\raggedright}{\thesection}{0.7em}{}
\usepackage{fancyhdr}
\pagestyle{fancy}\fancyhf{}
\fancyhead[L]{\small\sffamily Mathematical problem statements}
\fancyhead[R]{\small Source edition 22}
\fancyfoot[C]{\thepage}
\setlength{\parindent}{0pt}
\setlength{\parskip}{5pt plus 1pt}
\setlength{\emergencystretch}{3em}
\setcounter{tocdepth}{1}
\setcounter{secnumdepth}{1}
\setlist[itemize]{leftmargin=3em,itemsep=5pt,topsep=4pt,parsep=0pt}
\allowdisplaybreaks
\newcommand{\N}{\mathbb{N}}
\newcommand{\Z}{\mathbb{Z}}
\newcommand{\Q}{\mathbb{Q}}
\newcommand{\R}{\mathbb{R}}
\newcommand{\C}{\mathbb{C}}
\newcommand{\F}{\mathbb{F}}
\newcommand{\A}{\mathbb{A}}
\newcommand{\PP}{\mathbb P}
\newcommand{\E}{\mathbb{E}}
\newcommand{\eps}{\varepsilon}
\newcommand{\dd}{\,\mathrm d}
\newcommand{\sref}[2]{\hyperlink{source:#1:#2}{[S#2]}}
\newcommand{\eref}[2]{\hyperlink{extra:#1:#2}{[E#2]}}
% Turn the next switch off for a shorter reading copy. The quotations preserve
% every exactTarget field, including qualifications omitted from a short summary.
\newif\ifshowcatalogue
\showcataloguetrue
\newcommand{\cataloguescope}[1]{\ifshowcatalogue
 \paragraph{Exact scope in the supplied catalogue.}
 \begin{quote}\small #1\end{quote}\fi}

% Standalone wrapper; no external inputs or bibliography files are required.
\fancyhead[L]{\small\sffamily Research notebook: section 3433}
\fancyhead[R]{\small 27 September 2026}
\hypersetup{pdftitle={Research attempt 205},pdfauthor={Research notebook}}
\setlength{\parskip}{4pt plus 1pt}
\setlist[itemize]{before=\small\interlinepenalty10000,itemsep=3pt}
\begin{document}
\setcounter{section}{0}
% ================================================================
% problemId: 3433
% Canonical problem ID: 3433
% exactTarget SHA-256: 9abe39ba65eb8c8b840a3b27d4436223bee399af7be96c69046a532f71b7a3aa
\clearpage
\section*{\texorpdfstring{Jacob Palis Conjecture(Finitude of Attractors)(Dynamical Systems)}{Jacob Palis Conjecture(Finitude of Attractors)(Dynamical Systems)}}\label{3433}
\subsection{Definitions and mathematical statement}
For a diffeomorphism $f$ of a compact manifold, let
$B(A)=\{x:\operatorname{dist}(f^n(x),A)\to0\}$; for a flow replace $n$ by continuous positive time. In the selected Palis convention an attractor is a compact invariant transitive set with positive-volume basin, not necessarily an attractor defined by a trapping neighborhood. For each finite $r\ge1$, the target is density in both $\operatorname{Diff}^r(M)$ and the space of $C^r$ vector fields of systems with finitely many such sets $A_1,\ldots,A_k$ satisfying
\[
 \operatorname{vol}\left(M\setminus\bigcup_{i=1}^k B(A_i)\right)=0.
\]
The number $k$ may vary with the system. Both density assertions are required; physical measures and stochastic stability are not included in this selected part of the global conjecture.
\par\smallskip\textit{Formulation and concept references:} \sref{3433}{1}.
\cataloguescope{For every compact connected boundaryless smooth manifold M and every finite integer r >= 1, prove the conjunction of the following two density assertions: (i) in Diff\textasciicircum{}r(M) with its C\textasciicircum{}r topology there is a dense subset D whose elements f have finitely many attractors whose basins of attraction together cover M up to a set of smooth-volume measure zero; (ii) the same holds in the space X\textasciicircum{}r(M) of C\textasciicircum{}r vector fields with its C\textasciicircum{}r topology, for the attractors and basins of the generated flows. Attractor follows Palis (2005) Sections 2.1 and 2.3: a compact invariant set A that is transitive (has a dense orbit) and attracts a set of positive volume, its basin being the set of points whose forward orbit (or forward flow trajectory) converges to A. The finite number of attractors may depend on the element of D; no uniform bound is required. Volume-zero is independent of the choice of smooth positive density on M. This is the first bullet of Palis's Global Conjecture only; physical measures, stochastic stability, metric stability under parameter perturbation, and one-dimensional-family statements are separate conjectures not included.}
\subsection{Short English statement}
Can every smooth dynamical system be approximated by one in which almost every initial point is attracted to one of finitely many transitive attractors?
% BEGIN RESEARCH ADDITION 205
\subsection{Research attempt}

\paragraph{Current frontier.}
The incorporated definition has been checked in Palis, Sections 2.1 and
2.3, and the selected assertion is the \emph{first} bullet of Section 2.7.
It concerns distance convergence to compact transitive invariant sets;
the subsequent requirements about physical measures and stochastic
stability are not part of this target.\sref{3433}{1}
Berger's 2025 manuscript, published in the 2026 ICM proceedings, again
states the finite-attractor density question as Problem 1.4. Its discussion
of generic infinitely many sinks, and its stronger parameter-family
phenomena, distinguishes these from a disproof of density.\eref{3433}{1}
Hyperbolic and one-dimensional models furnish important solved settings,
but approximating every system by a uniformly hyperbolic one is not a
valid general strategy.\sref{3433}{1}\eref{3433}{1}

\paragraph{Attack route.}
A hyperbolic approximation route would require more than is available in
wild regions. A statistical route would have to upgrade convergence of
empirical measures to actual convergence in distance; that upgrade is
not automatic. The route investigated here is geometric: preserve a
\emph{fixed finite} collection of transitive trapping cores while enlarging
its finite-time capture set. An explicit quantitative persistence lemma
makes a convergent perturbation scheme rigorous \emph{provided} the basin
can be enlarged without adding indefinitely many cores. We then test both
that proviso and a tempting genericity strengthening against examples.

\paragraph{Attempt.}
Normalize a smooth positive volume $m$ so that $m(M)=1$. Changing this
normalization has no effect on null sets.

\textbf{1. Finite-time capture is a genuinely open condition.}
Let $f\in\operatorname{Diff}^r(M)$, and let $U_1,\ldots,U_k$ be nonempty
open sets with
\[
 f(\overline U_i)\subset U_i,\qquad
 A_i(f)=\bigcap_{n\ge0}f^n(\overline U_i).
 \tag{205.1}
\]
Here $M$ is compact, so the closures are compact. Suppose, as an explicit
additional hypothesis, that each $A_i(f)$ is transitive. Then
$U_i\subset B(A_i(f))$: the compact sets $f^n(\overline U_i)$ are nested,
and their intersection is $A_i(f)$, so their maximal distance to that
intersection tends to zero. The intersection is nonempty and invariant
under the homeomorphism $f$. Positive-volume basins follow since $U_i$
is nonempty and open. These are special cases of the catalogue's
attractors, not a redefinition of all attractors.

Set $U=\bigcup_iU_i$ and
\[
 C_N(f)=\bigcup_{j=0}^{N}f^{-j}(U),\qquad
 C_\infty(f)=\bigcup_{N\ge0}C_N(f).
\]
If $m(C_\infty(f))=1$, the selected finitely many basins cover almost
all of $M$. Conversely, a trapping-core argument only needs to prove
this capture property, not a mixing rate or a physical measure theorem.

\textbf{Lemma 205.1 (finite-capture persistence).}
For fixed open $U$, fixed finite $N$, and $\varepsilon>0$, the condition
$m(C_N(f))>1-\varepsilon$ is open in the $C^r$ topology.

\emph{Proof.}
By inner regularity choose compact $K\subset C_N(f)$ with
$m(K)>1-\varepsilon$. The continuous function
\[
 x\longmapsto \max_{0\le j\le N}
 \operatorname{dist}\bigl(f^j(x),M\setminus U\bigr)
\]
is positive on $K$ and has positive minimum $\delta$. Finite iterates
of maps $g$ sufficiently close to $f$ are uniformly within $\delta/2$
of the corresponding iterates of $f$. Consequently $K\subset C_N(g)$.
When $U=M$, the conclusion is immediate without the distance formula.
\hfill$\square$

There is an explicit estimate behind this openness. If $f$ is
$L$-Lipschitz and $\sup_xd(f(x),g(x))\le\eta$, then induction gives
\[
 \sup_xd(f^j(x),g^j(x))
 \le\eta\sum_{a=0}^{j-1}L^a. \tag{205.2}
\]
Thus the above compact capture certificate is retained whenever the
right side is below $\delta/2$ for $j\le N$. This estimate makes clear
that long capture times may force extremely small perturbations; no
uniform bound in $N$ has been smuggled into the argument.

\textbf{2. What a successful diagonal perturbation would actually require.}
Work in a local $C^r$ Banach chart and inside a complete closed ball
contained in the space of diffeomorphisms. Assume a neighborhood
$\mathcal W$ has fixed $U_1,\ldots,U_k$ for which (205.1) and transitivity
hold for every system under consideration, including any limiting system
in the chosen closed ball. This is a substantive persistence hypothesis.

Suppose one can choose nonempty nested closed balls $\mathcal V_n$
contained in that ball, with diameters tending to zero, and integers
$N_n$ such that
\[
 \mathcal V_{n+1}\subset\mathcal V_n,\qquad
 m(C_{N_n}(g))>1-2^{-n}
 \quad\text{for every }g\in\mathcal V_n. \tag{205.3}
\]
Completeness gives a unique $g_\infty\in\bigcap_n\mathcal V_n$.
For this \emph{one} map, (205.3) gives
$m(C_\infty(g_\infty))\ge1-2^{-n}$ for every $n$, hence full measure.
The same fixed $k$ transitive attractors from (205.1) therefore suffice.
This is a rigorous sufficient criterion; neither $k$ nor the times $N_n$
need be uniform across different initial neighborhoods.

Lemma 205.1 shows how a finite-time capture certificate at a newly chosen
map can be protected in the next ball. It does not construct that map.
The missing dynamical step is the following pathwise basin-completion
property: inside a sequence of increasingly constrained neighborhoods,
can one improve the captured measure arbitrarily close to one while
retaining a finite set of trapping cores and its transitivity? The
question is about the same evolving construction, not independent
approximations with arbitrarily many attractors.

For vector fields take blocks satisfying
$\varphi_t^X(\overline U_i)\subset\overline U_i$ for $t\ge0$ and
$\varphi_t^X(\overline U_i)\subset U_i$ for $t>0$, and set
\[
 A_i(X)=\bigcap_{t\ge0}\varphi_t^X(\overline U_i),\qquad
 C_T(X)=\{x:\varphi_t^X(x)\in U\text{ for some }0\le t\le T\}.
\]
Forward trapping for all $t\ge0$ makes these compact sets nested;
the intersection is invariant under the full flow. On a compact phase
space, finite-time continuous dependence on the vector field proves the
analogue of Lemma 205.1 by the same compact-certificate argument. Nested
balls in $\mathcal X^r(M)$ then give the analogue of (205.3). One must
still supply the basin-completion step for vector fields separately.
A perturbation of a time-one diffeomorphism need not be a time-one map
of a nearby vector field, so it cannot establish this second branch.

\textbf{3. An explicit smooth failure of ``arbitrarily nearly full'' capture.}
The need to control the number of cores can be tested without wild
higher-dimensional dynamics. In a coordinate near $0\in S^1$, take a
smooth vector field $X=v(x)\,\partial_x$ with
\[
 v(x)=
 \begin{cases}
  e^{-1/x^2}\sin(1/x),&x>0,\\
  0,&x=0,\\
  e^{-1/x^2},&x<0.
 \end{cases} \tag{205.4}
\]
Use (205.4) in a sufficiently small coordinate interval; choose its
positive endpoint inside a positive lobe and extend $v$ positively
around the remaining arc of the circle. Such an extension is smooth:
one glues within positive collars, and all derivatives at $0$ vanish.
It introduces no additional zeros.

For all sufficiently large $n$, the positive zeros $x_n=1/(n\pi)$
accumulating at zero satisfy
\[
 v'(x_n)=-(-1)^n n^2\pi^2 e^{-n^2\pi^2}.
\]
For even $n$ these are hyperbolic sinks; for odd $n$ they are sources.
Every point in an interval between adjacent zeros flows monotonically
toward its sink endpoint. Points in the complementary arc flow toward
$0$ from the negative side. Thus the countably many sink singletons,
together with $\{0\}$, have basins covering all points except a countable
set of sources. Each is transitive and has positive-volume basin.

By continuity from below of measure, for every $\varepsilon>0$ some
\emph{finite} subcollection has basin union of measure exceeding
$1-\varepsilon$. Nevertheless no finite collection of transitive
attractors can have a full-measure basin union for this particular field.
Here is a proof that does not assume attractors are minimal.
If $x$ tends to a sink $s$ and also belongs to $B(A)$ for a compact
invariant set $A$, then $s\in A$. If the basins of finitely many $A_i$
covered almost all of $M$, every positive-volume sink basin would force
its sink to belong to some $A_i$. But a transitive set cannot contain
two distinct sinks: their disjoint forward-invariant trapping
neighborhoods would both have to be visited by its dense orbit, and
once the orbit enters the first it cannot visit the second. This
argument works even if ``dense orbit'' is interpreted as a two-sided
orbit, since the two visit times can still be ordered. Infinitely many
sinks therefore cannot be assigned to finitely many such $A_i$.

The time-one map of (205.4) is a $C^\infty$ circle diffeomorphism with
the same forward limit points and the same obstruction. Hence this
failure of the approximate-coverage inference occurs in both dynamical
categories in the question. It is \emph{not} a counterexample to either
density assertion: a single bad system is allowed.

Indeed, the vector-field example is especially easy to approximate by
good systems. More generally, every $C^r$ vector field on $S^1$ can first
be approximated in $C^r$ by a smooth one $v(x)\partial_x$. Choose an
arbitrarily small constant $c$ such that $-c$ is a regular value of $v$.
Sard's theorem in dimension one supplies such constants. Then $v+c$
either has no zeros, in which case its whole-circle periodic orbit is a
transitive attractor, or has finitely many simple zeros. In the latter
case the zeros alternate between sinks and sources, and all points
except the sources lie in the finitely many sink basins. This proves
the vector-field branch on the circle by a direct classical argument.
It does not prove the higher-dimensional assertions or the general
$d$-dimensional diffeomorphism branch.

\textbf{4. Why replacing geometry by time averages also loses a hypothesis.}
There is a short exact model of the obstruction. In the two-sided shift
space $\{0,1\}^{\Z}$, take $x$ with symbol $1$ at positive coordinates
$2^k$ and $0$ elsewhere. For every fixed window, the set of iterates
whose window contains a $1$ has asymptotic density zero. Approximation
of continuous functions by finite-window functions therefore gives
\[
 \frac1N\sum_{j=0}^{N-1}\delta_{\sigma^j x}
 \ \Longrightarrow\ \delta_{0^{\Z}}.
\]
Yet at times $j=2^k$ the central symbol is $1$, so
$d(\sigma^j x,0^{\Z})$ is bounded below in a standard product metric.
Distance convergence fails. This is a logical countertest for the
statistical-to-geometric inference, not a manifold counterexample or
an assertion about a positive-volume family of shift points.

\textbf{5. A genericity strengthening of the repair scheme is impossible.}
Suppose one strengthened step 2 by claiming, throughout a nonempty open
set $\mathcal W$ with fixed persistent transitive trapping cores, that
for every $\varepsilon>0$ the finite-time capture conditions were dense.
They are open by Lemma 205.1. The Baire theorem would then make full
capture by these finitely many cores a residual property in
$\mathcal W$. In a Newhouse open set with residual infinitely many sinks,
this contradicts the two-sink argument above. Such open sets occur in
the settings described in Berger's Section 1.2.\eref{3433}{1}
Thus an unrestricted \emph{open-dense} basin-completion upgrade is not
an available missing lemma. A successful density construction in those
regions must be able to select an exceptional, potentially meagre set
of systems rather than prove a generic finite-core theorem there.

\paragraph{Stress test.}
The compact-certificate argument avoids unjustified interchange of
limits of systems and limits of basins. It first fixes a finite time,
protects the resulting compact set under perturbation, and only then
takes a limit of systems. The finite family and transitivity are explicit
hypotheses and have not been inferred from compactness or chain
recurrence. The example (205.4) shows why allowing that family to grow
without bound is not harmless, even for smooth flows.

The argument also separates three quantifiers that cannot be exchanged:
one map for all accuracy levels, a possibly different finite collection
for each accuracy level, and a dense set of maps versus a residual set
of maps. The Newhouse stress test rejects a stronger proposed route,
not Palis's density assertion. Smooth-volume and Baire-category
statements have never been identified.

\Needspace{10\baselineskip}
\paragraph{Outcome.}
\textbf{Outcome: Exploratory approach with unresolved gap.}

The attempt proves finite-time capture persistence with an explicit
perturbation estimate and gives a valid pathwise completion criterion
for both maps and flows. A smooth circle construction disproves the
inference from arbitrarily large finite basin coverage to exact finite
full coverage; a separate shift example disproves the inference from
statistical convergence to distance convergence. The vector-field case
on $S^1$ is recovered, without a claim of novelty.

No basin-completing perturbation has been constructed for arbitrary
systems while keeping finitely many transitive cores. Moreover, the
most convenient generic version of that missing step is incompatible
with Newhouse phenomena. Resolving the selected conjunction would need
an exceptional-system construction or another mechanism that works
for every finite $r$ in both categories without silently imposing
trapping neighborhoods, physical measures, or hyperbolicity on the
original target.

% Keep the sources heading with a useful initial bibliography block.
\Needspace{8\baselineskip}
% END RESEARCH ADDITION 205
\subsection{Sources}
\begin{itemize}\raggedright
\item[\textnormal{[S1]}]\hypertarget{source:3433:1}{} J. Palis, A global perspective for non-conservative dynamics, Ann. IHP AN 22 (2005), 485-507.
\url{https://ems.press/content/serial-article-files/15975}.
{\small\textit{Catalogue formulation source.}}
% BEGIN REFERENCE ADDITION 205
\item[\textnormal{[E1]}]\hypertarget{extra:3433:1}{}
P. Berger, \emph{Wild dynamics on manifolds}, arXiv:2510.12929v1
(14 October 2025), Sections 1.1--1.2, especially Problem 1.4 and
Theorem 1.7; published in the Proceedings of the International Congress
of Mathematicians 2026, pp. 24--41, DOI: 10.1137/25M1810787.
The genericity statements are used with their stated topology and
model restrictions, not as counterexamples to density.
\url{https://arxiv.org/abs/2510.12929}.
\url{https://epubs.siam.org/doi/10.1137/25M1810787}.
% END REFERENCE ADDITION 205
\end{itemize}

\end{document}
